\documentclass[book.tex]{subfiles}
\begin{document}

\chapter{Clifford Group Equivariant Networks}

\label{chap:deep_onet}
\noindent\rule{11cm}{0.4pt} \\
\textbf{Prerequisites:} \autoref{chap:dft} \\
\textbf{Difficulty Level:} ** \\
\noindent\rule{11cm}{0.4pt}
\vspace{5mm}

% https://arxiv.org/abs/2305.11141
% https://github.com/DavidRuhe/clifford-group-equivariant-neural-networks

%Abstract: In this talk, I will present Clifford Group Equivariant Neural Networks, an innovative method for building E(n)-equivariant networks based on Clifford (geometric) algebras. First, I will give an introduction to the Clifford algebra and its geometric applications. Then, I will introduce the Clifford group and how it always acts through the orthogonal group. As such, a parameterization that is equivariant to the Clifford group will automatically be equivariant to the orthogonal group of, e.g., rotations and reflections. We show that any polynomial (under the algebra’s geometric product) is such a parameterization.  We propose several layers from these insights and conduct experiments in three-, four-, and five-dimensional spaces. One of these experiments even includes equivariance to the nondefinite O(1,3) Lorentz group from the same code implementation. Finally, I will provide guidance on how to utilize our codebase for implementing these algorithms.

%\url{https://micde.umich.edu/event/sciml-webinar-david-ruhe-geometric-clifford-algebra-networks/}

\section{Introduction}

Most people have focused on $SE(2)$ or $SE(3)$ or $SO(3)$. What kind of algebraic tools can we use to enable neural networks to learn geometry?

Consider a graph structured equivariant structure. We can categorize such networks a few categories
\begin{itemize}
    \item Regular representations: We use a group convolution. This is most famously done by the convolutional neural network. We perform an integral over the group, but this is challenging for Lie Group such as continuous groups.
    \item Scalarization methods: Operate with invariant features unaffected by transformations. Or update vector-valued information using scalar multiplication. This limits expressibility
    \item Steerable networks: Take tensor products of features. Representations of group acting on these is Wigner-D matrices. You have to decompose into irreps using Clebsch-Gordon coefficients.
\end{itemize}

This chapter looks at the Clifford algebra, or the geometric algebra. These objects consist of scalar, vectors, bivectors, and trivectors more generally called multivectors. We allow neural networks to operate on multivectors.

\section{The Clifford Algebra}

We first start with bilinear forms. We need a notion of distance provided by the bilinear form
\begin{align}
    \langle \cdot, \cdot \rangle \mapsto F
\end{align}
Can use inner product to compute angles and distances.

This introduces symmetry from the orthogonal group (the group of linear transformations preserving this bilinear form is the orthogonal group). For example, consider reflections and rotations. 

An Euclidean inner product is given by
\begin{align}
    v^TUw
\end{align}
where $U$ is a PST matrix (?). This framework works for non-Euclidean metrics as well. The orthogonal group of the Lorentz group is $O(1, 3)$. However the same definition holds.

The Clifford algebra is also called the geometric algebra. An algebra is a vector space equipped with a product. 

\begin{align}
    uv \in R, u \in R, v \in R
\end{align}
To give this expression meaning, we have to connect this product to the geometry. Inner products give scalars. We want an inner product not restricted to scalars. Here are axioms
\begin{align}
    v^2 &= \langle v, v \rangle \\
    (u + v)^2 &= uv + vu + uv + vu 
\end{align}
The last condition holds if and only if
\begin{align}
    uv + vu = 2\langle u, v\rangle
\end{align}
We still need to define the space in which this individual product lives.
\begin{align}
    uv &= \frac{uv + vu}{2} + \frac{uv-vu}{2}  \\
    &= \langle u, v \rangle + u \wedge v
\end{align}
Here $\wedge$ is the wedge product, an anticommutative product. It computes the area of the parallelogram the two vectors define.

\begin{figure}
    \centering
    % \includegraphics[]{figures/Differentiable Physics/Neural Diff Eq/Clifford Algebra/wedge_product.png}
    \includegraphics[]{figures/Differentiable Physics/Neural Diff Eq/Clifford Algebra/wedge_product1.png}
\end{figure}
Can we find $w$ such that
\begin{align}
    v \cdot w &= 0
\end{align}
and that the area with $v$ is the same as $u \wedge v$
\begin{align}
    wv &= u \wedge v = \frac{1}{2}(uv - vu) \\
    w &= \frac{u - vuv^{-1}}{2}
\end{align}

\subsection{The Algebra Basis}

Start with $\mathbb{R}^3$ with basis $e_1$, $e_2$, $e_3$. Then we have
\begin{align}
    e_1 e_2 &= \langle e_1, e_2 \rangle + e_1 \wedge e_2 \\
    &= e_1 \wedge e_2
\end{align}
since the basis vectors are orthogonal. So in this case, it equals the wedge product. This new object is called a basis plate. We then create the
\begin{align}
    \mathrm{Cl}(\mathbb{R}^3, \langle \cdot, \cdot \rangle )
\end{align}
\begin{figure}
    \centering
   
     % \includegraphics[ht]{figures/Differentiable Physics/Neural Diff Eq/Clifford Algebra/basis.png}
     \includegraphics[ht]{figures/Differentiable Physics/Neural Diff Eq/Clifford Algebra/basis1.png}
\end{figure}
For orthogonal basis vectors we have
\begin{align}
    e_1e_2 &= -e_2 e_1
\end{align}

Any multivector can be written in a basis of scalars, vectors, bivectors, and trivectors. 
\begin{align}
    x &= x_0 1 + x_1 e_1 + x_2 e_2 + x_3 e_3 + x_{12} e_{12} + x_{13}e_{13} + x_{23} e_{23} + x_{123} e_{123}
\end{align}
A vector in the algebra is an element of the algebra with non-vector parts of this basis set to 0.
\subsection{The Geometric Product}
Let us take
\begin{align}
    a, b \in \mathrm{Cl}(\mathbb{R}^2, q),. q(e_i) = \langle e_i, e_1 \rangle = 1 
\end{align}
Let us take the product
\begin{align}
    ab &= (a_1 e_1 + a_2 e_2)(b_1 e_1 + b_2 e_2) \\
    & a_1e_1(b_1 e_1 + b_2 e_2) + a_2 e_2 (b_1 e_1 + b_2 e_2) \\
    &= a_1 b_1 e_1^2 + a_1 b_2 e_1 e_2 + a_2 b_1 e_2 e_1 + a_2 b_2 e_2^2 \\
    &= (a_1 b_1 + a_2 b_2) + \cdots
\end{align}

\subsection{Clifford Algebra Representations}

The Clifford group is given by
\begin{align}
\Gamma^{[X]}
\end{align}
This is a subset of the Clifford algebra. Its action $\rho(w)$ is the twisted conjugation. We consider reflections in particular. Every orthogonal transformation decomposes into multiple reflections in the Clifford group. This action satisfies
\begin{align}
    \rho(w)(x_1 + x_2) &= \rho(w)(x_1) + \rho(w)(x_2) \\
    \rho(w)(x_1 x_2) &= \rho(w)(x_1) \rho(w)(x_2) \\
    \rho(w)(\phi x) &= \phi \rho(w)(x) \\
    \langle \rho(w)(x_1), \rho(w)(x_2) \rangle &= \langle x_1, x_2 \rangle
\end{align}
This $\rho$ is a representation of the orthogonal group. 
\subsection{Polynomials and Grade Projections}

Using the properties of $\rho(w)$ we can find two general operations. First consider polynomials,
\begin{align}
F: \mathrm{Cl}(V, q)^\ell \to \mathrm{Cl}(V, q)
\end{align}
Then we have
\begin{align}
    \rho(w)(F(x_1, \dotsc, x_\ell)) &= F(\rho(w)(x_1),\dotsc, \rho(w)(x_\ell))
\end{align}
As an example consider
\begin{align}
    F(x_1, x_2) &= \phi_0 + \phi_1 x_1 + \phi_2 x_2 + (\phi_{11}x_1^2 + \phi_{12} x_{12} + \phi_{21} x_{21} + \phi_{22} x_2^2)
\end{align}
Let 
\begin{align}
    (_)^m: \mathrm{Cl}(V, q) \to \mathrm{Cl}(V, q)
\end{align}
denote the projection onto the grade $m$ par
\begin{align}
    (_)^1(\phi_0 + \phi_1 x_1 + \phi_2 x_2 + (\phi_{11}x_1^2 + \phi_{12} x_{12} + \phi_{21} x_{21} + \phi_{22} x_2^2)) &= \phi_1 x_1 + \phi_2 x_2
\end{align}
These grade projections are equivariant

\section{Equivariant Layers}
\subsection{Linear Layers}
A polynomial up to the first order is just a linear map. Let us have
\begin{align}
    x_1, \dotsc, x_\ell
\end{align}
View these as channels. We can linearly combine
\begin{align}
    T_{\phi, j} &= \sum_{i} \phi_{ji}x_i
\end{align}
where $\phi_{ji} \in \mathbb{R}$. 
\subsection{Parameterized Geometric Product}
When taking a product of two Clifford alebra elements, you can decompose in terms of grades
\begin{align}
P_\phi(x_1, x_2)^{(k)} &= \sum_{i} \sum_j \phi_{ijk} (x_1^{(i)} x_2^{(j)})^{(k)} \dotsc
\end{align}
Each of these grade outputs can be viewed as a different geometric product. Here again we have
\begin{align}
    \phi_{ijk} \in \mathbb{R}
\end{align}

\section{Network Architecture}
You can combine linear layers, parameterized geometric products. You can also use gated nonlinearities as used elsewhere in the equivariant literature. This architecture can be tested on the $n$-body dataset. In the three dimensional case, this system is the same as the SEGNN.

This architecture can also be applied for $O(1, 3)$ jet-tagging for top-tagging at CERN. Top-tagging is the problem of taking collision data (momenta, energy) of hundreds of particles and giving a binary classification of whether a top quark was produced. These particles are moving at relativistic speeds, so we must respect Lorentz equivariance. By using the infrastructure above and swapping the metric for the $O(1, 3)$ metric, you can achieve a Lorentz-equivariant network. This achieves performance similiar to LorentzNet.

\section{Conclusion}
A powerful advantage of this framework is avoiding the need for the group convolution. It also achieves state of the art results for different groups. The models have computational overhead similar to other steerable networks though.
\end{document}
